%---------------------------------------------------------------
% verwendete packages
%---------------------------------------------------------------
\usepackage[english]{babel}
% neue deutsche Trennmuster usw.
\usepackage[utf8]{inputenc}       % Quelltext ist Latin-1 (d.h. Windows-) kodiert (direkte Eingabe von Umlauten usw.)
\usepackage{longtable}              % Tabellen über mehrere Seiten
\usepackage[T1]{fontenc}            % Cork-encoding für erweiterte Zeichensätze (entweder mit ae oder cm-super-fonts verwenden)
\usepackage[scale=0.92]{helvet}
\usepackage{textcomp}               % erweiterte Sonderzeichen
\usepackage{fancyhdr}               % beliebige Kopfzeilendefinition
\usepackage{graphicx}               % Einbinden von Grafiken
\usepackage{amssymb,amsbsy,amsmath} % AMS
\renewcommand{\familydefault}{\sfdefault}
%\usepackage[slantedGreek]{mathpazo} % Schriftart Palatino auswählen
\usepackage{eso-pic}                % Für Wasserzeichen für Draft
\usepackage{verbatim}               % Verbatimumgebung
\usepackage{fancyvrb}               % Verbatimumgebung
\usepackage{color}                  % Farbdefinitionen
\usepackage{xspace}                 % vermeidet Fehler von Leerzeichen nach eigenen Befehlen
\usepackage[rel]{overpic}           % Bilder bearbeiten
\usepackage{colortbl}               % Farbe in Tabellen
\usepackage{textcomp}               % Sonderzeichen
\usepackage{floatflt}               % umflossene Grafiken
\usepackage{array}                  % erweiterte Felder
\usepackage{booktabs}               % verbessertes Tabellenlayout
\usepackage{float} % in the preamble
\usepackage{csquotes}
\usepackage{makecell} % add this in your preamble
\usepackage{listings}
\usepackage{xcolor}
\usepackage{caption}
\usepackage{qrcode}

\usepackage{pdfpages}


\lstset{ backgroundcolor=\color{gray!10}, basicstyle=\ttfamily\small, frame=single, breaklines=true, numbers=left, numberstyle=\tiny\color{gray}, keywordstyle=\color{blue}, commentstyle=\color{gray}, stringstyle=\color{orange}, }

\lstdefinelanguage{json}{
    basicstyle=\ttfamily\small,
    numbers=left,
    numberstyle=\tiny\color{gray},
    stepnumber=1,
    numbersep=5pt,
    showstringspaces=false,
    breaklines=true,
    frame=lines,
    literate=
     *{:}{{{\color{black}:}}}{1}
      {,}{{{\color{black},}}}{1}
      {\{}{{{\color{black}\{}}}{1}
      {\}}{{{\color{black}\}}}}{1}
      {[}{{{\color{black}[}}}{1}
      {]}{{{\color{black}]}}}{1},
}
%
%---------------------------------------------------------------
% nicht verwendete packages
%---------------------------------------------------------------
%\usepackage{multicol}               %
%\usepackage{ae}                     % ae-Zeichensätze für pdf-Darstellung, aber Umlaute nicht gut (Trennung!)
%\usepackage{mathptmx}               % Times verwenden
%\usepackage[scaled]{helvet}         % Helvetica
%\renewcommand{\familydefault}{\sfdefault} % auf serifenlose Schrift umschalten
%\usepackage{cmbright}               % cm bright
%\usepackage{courier}
%\usepackage{pifont}
%\usepackage{showkeys} % im dvi-previewer werden verweise angezeigt
%\usepackage{Styles/Sonstige/ifmfont}% Frutiger verwenden
%\usepackage{eqmark}                 % Formellayout
%\usepackage{framed}                 % umrandeter Text
%\usepackage{lscape}                 % Querformat
%\usepackage{pdflscape}              % Querfomrat für Reader automatisch drehen
%\usepackage{ifthen}                 % zum Prüfen, ob Packages eingebunden
%\usepackage{longtable}              % Tabellen über mehrere Seiten
%\usepackage{natbib}                 % Style Bibliography
%\usepackage[figuresright]{rotating} % Gleitumgebungen drehen (beidseitig geht nicht richtig... :-( )
%\usepackage{textcomp}               % Style Bibliography
%\usepackage[ngerman,USenglish]{babel}% deutsche Trennmuster für References ermöglichen
%
%---------------------------------------------------------------
% Mehrere Literaturverzeichnisse aus einer *.bib-Datenbank
%---------------------------------------------------------------
%\usepackage{multibib}
%\newcites{eigene}{ }
%\newcites{andere}{ }
%---------------------------------------------------------------
% 
\newfont{\fonttimes}{ptmr at 12pt}
\newfont{\fonthelvet}{phvr at 11pt}
\newfont{\fontpalatino}{pplr at 12pt}
\newfont{\fontcm}{cmr12}
